% !TEX root = ../main.tex
\section{Theoretical Framework}

This section summarises the physical ingredients required by \texttt{tpeanuts}. The objective is not to provide a complete review of neutrino physics, but to establish the conventions that must be preserved when updating PEANUTS to a tensor-based implementation.

\subsection{Neutrino propagation fundamentals}
\label{sec:propagation-fundamentals}

This subsection collects the formalism common to every propagation regime considered in this thesis: the definition of the mixing matrix, vacuum oscillations, and the modification of the propagation Hamiltonian by ordinary matter. The specific neutrino sources and new-physics extensions built on top of this common core are treated separately in the following subsections.

\subsubsection{Neutrino mixing}

One of the fundamental discoveries in particle physics is that neutrinos are produced and detected in flavour eigenstates, whereas they propagate through space as mass eigenstates. Since these two bases do not coincide, a neutrino created with a definite flavour evolves into a quantum superposition of different flavours during propagation, giving rise to the phenomenon of neutrino oscillations \parencite{Giunti2007}.

Within the Standard Model extended to include neutrino masses, the flavour eigenstates \(\ket{\nu_\alpha}\), \(\alpha=e,\mu,\tau\), are related to the mass eigenstates \(\ket{\nu_i}\), \(i=1,2,3\), through a unitary transformation known as the Pontecorvo--Maki--Nakagawa--Sakata (PMNS) matrix \parencite{Pontecorvo1957,Pontecorvo1958,MakiNakagawaSakata1962},

\begin{equation}
  \ket{\nu_\alpha} = \sum_{i=1}^{3} U_{\alpha i}^{*}\ket{\nu_i},
\end{equation}

or, equivalently,

\begin{equation}
  \ket{\nu_i} = \sum_{\alpha=e,\mu,\tau} U_{\alpha i}\ket{\nu_\alpha}.
\end{equation}

The matrix \(U\) is unitary, \(U^{\dagger}U=UU^{\dagger}=\mathbb{I}\), which guarantees the conservation of probability during neutrino propagation.

\paragraph{Physical interpretation.}
The matrix elements \(U_{\alpha i}\) represent the probability amplitudes that a neutrino with flavour \(\alpha\) is found in the mass eigenstate \(i\). Consequently, \(|U_{\alpha i}|^2\) gives the probability that a neutrino produced as flavour \(\alpha\) propagates in the mass eigenstate \(i\). This mismatch between flavour and mass eigenstates is directly analogous to the quark sector, where flavour mixing is described by the Cabibbo--Kobayashi--Maskawa (CKM) matrix \parencite{PDG2022}. In the lepton sector, however, the mixing angles are significantly larger, making flavour oscillations observable over laboratory, atmospheric, solar and astrophysical baselines.

\paragraph{Standard parametrisation.}
For three active Dirac neutrinos, the PMNS matrix is written as the product of three elementary rotations, \(U_{\mathrm{PMNS}}=R_{23}(\theta_{23})\,U_{13}(\theta_{13},\delta_{\mathrm{CP}})\,R_{12}(\theta_{12})\), yielding the standard Particle Data Group (PDG) parametrisation \parencite{PDG2022}

\begin{equation}
  U_{\mathrm{PMNS}} =
  \begin{pmatrix}
    c_{12}c_{13} & s_{12}c_{13} & s_{13}e^{-i\delta_{\mathrm{CP}}}\\
    -s_{12}c_{23}-c_{12}s_{23}s_{13}e^{i\delta_{\mathrm{CP}}} & c_{12}c_{23}-s_{12}s_{23}s_{13}e^{i\delta_{\mathrm{CP}}} & s_{23}c_{13}\\
    s_{12}s_{23}-c_{12}c_{23}s_{13}e^{i\delta_{\mathrm{CP}}} & -c_{12}s_{23}-s_{12}c_{23}s_{13}e^{i\delta_{\mathrm{CP}}} & c_{23}c_{13}
  \end{pmatrix},
\end{equation}

with \(c_{ij}=\cos\theta_{ij}\), \(s_{ij}=\sin\theta_{ij}\). Each rotation has a clear physical meaning: \(\theta_{12}\) governs mixing between the first and second mass eigenstates and dominates solar-neutrino oscillations together with the smaller mass splitting \(\Delta m_{21}^2\); \(\theta_{23}\) controls mixing between the second and third eigenstates and is primarily responsible for atmospheric-neutrino oscillations; \(\theta_{13}\) couples the first and third eigenstates and, although considerably smaller than the other mixing angles, plays a crucial role in matter effects and enables sensitivity to CP violation. The complex phase \(\delta_{\mathrm{CP}}\) introduces a difference between neutrino and antineutrino oscillation probabilities, providing a possible source of CP violation in the lepton sector.

\paragraph{Oscillation parameters.}
Oscillation experiments are sensitive not to the absolute neutrino masses but to the differences between their squared masses, \(\Delta m_{ij}^2 = m_i^2 - m_j^2\). Only two independent mass-squared differences exist, \(\Delta m_{21}^2\) and \(\Delta m_{31}^2\) (or equivalently \(\Delta m_{32}^2\)), together with the three mixing angles and the CP-violating phase,

\begin{equation}
  \{\theta_{12},\theta_{13},\theta_{23},\Delta m_{21}^2,\Delta m_{31}^2,\delta_{\mathrm{CP}}\}.
\end{equation}

These six quantities completely determine neutrino oscillation probabilities within the standard three-neutrino oscillation framework, with current best-fit values obtained through global analyses combining solar, atmospheric, reactor and accelerator data. The NuFIT collaboration's NuFIT~5.2 release \parencite{NuFit52} fixes the oscillation-parameter presets used throughout this work's forward-simulation and validation notebooks; the more recent NuFIT~6.0 update \parencite{EstebanNuFit2024} is used separately, only as an external benchmark in the real-data detector-level fits of \cref{sec:detector-validation-methodology}, not as a simulation input.

\subsubsection{Vacuum oscillations}

Once produced in a flavour eigenstate, neutrinos propagate in vacuum as coherent superpositions of mass eigenstates. Since each mass eigenstate accumulates a different quantum phase, their interference is the physical origin of neutrino oscillations \parencite{Giunti2007}. The evolution is governed by the Schr\"odinger equation,

\begin{equation}
  i\frac{d}{dt}\ket{\nu(t)} = H\ket{\nu(t)},
\end{equation}

where \(H\) is the effective Hamiltonian describing neutrino propagation.

\paragraph{Vacuum Hamiltonian.}
In vacuum, \(H\) is diagonal in the mass basis, \(H_{\mathrm{mass}}=\operatorname{diag}(E_1,E_2,E_3)\), with \(E_i=\sqrt{p^2+m_i^2}\). Since neutrinos produced in laboratory, solar and atmospheric experiments are highly relativistic, \(m_i^2\ll p^2\), the energy can be expanded as \(E_i\simeq p + m_i^2/(2E)\), with \(E\simeq p\) the neutrino energy. The common momentum term contributes only an overall global phase, \(e^{-ipL}\), which has no observable consequences and can be dropped, leaving the effective vacuum Hamiltonian

\begin{equation}
  H_{\mathrm{vac}} = \frac{1}{2E}\operatorname{diag}\left(m_1^2,m_2^2,m_3^2\right).
\end{equation}

Flavour evolution requires transforming \(H_{\mathrm{vac}}\) into the flavour basis through the PMNS matrix,

\begin{equation}
  H_{f} = U H_{\mathrm{vac}} U^{\dagger} = \frac{1}{2E}\,U\operatorname{diag}(m_i^2)U^{\dagger}.
\end{equation}

\paragraph{Evolution operator.}
For a constant Hamiltonian, the Schr\"odinger equation is solved by \(\ket{\nu(L)}=e^{-iHL}\ket{\nu(0)}\). For a neutrino produced in flavour \(\alpha\), \(\ket{\nu(0)}=\ket{\nu_\alpha}\), the transition amplitude to flavour \(\beta\) after a distance \(L\) is

\begin{equation}
  A_{\alpha\to\beta}(L,E) =
  \sum_{i=1}^{3} U_{\beta i}
  \exp\left(-i\frac{m_i^2 L}{2E}\right)
  U_{\alpha i}^{*},
  \label{eq:vacuum-transition-amplitude}
\end{equation}

with oscillation probability \(P_{\alpha\to\beta}=|A_{\alpha\to\beta}|^2\). Unitarity of the PMNS matrix guarantees \(\sum_\beta P_{\alpha\to\beta}=1\), i.e.\ conservation of probability.

\paragraph{Oscillation phases.}
Oscillations arise from the phase differences accumulated by the mass eigenstates, controlled by the mass-squared differences \(\Delta m_{ij}^2=m_i^2-m_j^2\), which define the phase entering the transition amplitude of \cref{eq:vacuum-transition-amplitude}, \(\Delta_{ij}=\Delta m_{ij}^2 L/2E\). Squaring a coherent two-state superposition introduces a half-angle: the two-flavour reduction of \cref{eq:vacuum-transition-amplitude} gives \(P=1-\sin^2(2\theta)\sin^2(\Delta_{ij}/2)\), so it is \(\Delta_{ij}/2\), not \(\Delta_{ij}\) itself, that governs the observable probability. In the practical units customary in oscillation phenomenology \parencite{PDG2022}, \(\Delta_{ij}/2=1.267\,\Delta m_{ij}^2\,[\mathrm{eV}^2]\,L\,[\mathrm{km}]/E\,[\mathrm{GeV}]\); oscillations are thus governed by the ratio \(L/E\) rather than by the baseline or the energy separately.

\subsubsection{Matter effects}

The vacuum formalism above assumes that neutrinos propagate through empty space. In many settings relevant to \texttt{tpeanuts}, such as the solar interior, the Earth's mantle and core, or a supernova envelope, neutrinos instead traverse dense media, where coherent forward scattering off the particles of the medium modifies the propagation Hamiltonian. This is the matter effect, also known as the Mikheyev--Smirnov--Wolfenstein (MSW) effect \parencite{Wolfenstein1978,MikheyevSmirnov1985,Giunti2007}. Because coherent forward scattering transfers essentially no momentum to the medium, it does not remove neutrinos from the beam; instead, it modifies the phase accumulated during propagation, which is described as an additional effective potential in the Schr\"odinger equation.

\paragraph{Matter potential.}
Neutrinos interact with matter through both neutral-current (NC) and charged-current (CC) weak processes. NC interactions, \(\nu_\alpha+f\to\nu_\alpha+f\) with \(f\) an electron, proton or neutron, affect all active flavours equally; since oscillation probabilities depend only on relative phases between flavours, this common contribution is unobservable and is omitted from the oscillation Hamiltonian. Electron neutrinos additionally undergo the CC process \(\nu_e+e^-\to\nu_e+e^-\), mediated by \(W^{\pm}\) exchange, which introduces a flavour-dependent potential

\begin{equation}
  V_{\mathrm{CC}} = \sqrt{2}\,G_F n_e,
\end{equation}

where \(G_F\) is the Fermi coupling constant and \(n_e\) the local electron number density. This additional, electron-flavour-specific interaction breaks the flavour symmetry of vacuum propagation and is the origin of the MSW effect.

\paragraph{Effective Hamiltonian.}
The full Hamiltonian governing flavour evolution in matter is obtained by adding the matter potential to the vacuum Hamiltonian,

\begin{equation}
  H = H_{\mathrm{vac}} + V_{\mathrm{mat}},
  \qquad
  H_{\mathrm{vac}} = \frac{1}{2E}\,U\operatorname{diag}\!\left(m_1^2,m_2^2,m_3^2\right)U^{\dagger},
  \qquad
  V_{\mathrm{mat}} =
  \sqrt{2}G_F n_e
  \begin{pmatrix}
    1 & 0 & 0\\
    0 & 0 & 0\\
    0 & 0 & 0
  \end{pmatrix},
\end{equation}

where \(V_{\mathrm{mat}}\) is diagonal in the flavour basis, since only the electron flavour experiences the CC interaction. For antineutrinos, both the matter potential changes sign, \(V_{\mathrm{mat}}\to-V_{\mathrm{mat}}\), and the mixing matrix is complex-conjugated, \(U\to U^{*}\) (equivalently \(\delta_{\mathrm{CP}}\to-\delta_{\mathrm{CP}}\)), a convention detailed in \cref{sec:antinu-convention}. In matter, the propagation basis becomes density-dependent, \(H=U_m\operatorname{diag}(\lambda_1,\lambda_2,\lambda_3)U_m^\dagger\), with \(U_m\) the effective mixing matrix and \(\lambda_i\) the effective eigenvalues; as the electron density varies along the trajectory, both evolve continuously, responsible for the characteristic energy dependence of solar-neutrino oscillations.

\subsection{Neutrino sources}

The oscillation formalism introduced above is universal: the same Schr\"odinger equation governs flavour evolution regardless of how the neutrino is produced. What varies between sources is the physical environment of production and, consequently, the density profile traversed before detection, which determines the relative importance of vacuum oscillations, matter effects and decoherence. Each source is therefore sensitive to a different region of the oscillation parameter space.

The current version of \texttt{tpeanuts} is designed primarily for the simulation of solar and atmospheric neutrinos, which together probe a broad range of energies, baselines and matter densities. Owing to its modular architecture, however, the propagation engine itself is independent of the neutrino source, so that additional sources can be incorporated with minimal modifications, as discussed for reactor and accelerator neutrinos below.

\subsubsection{Solar neutrinos}

Solar neutrinos are produced by thermonuclear fusion in the solar core, where temperatures reach \(\sim 1.5\times10^{7}\,\mathrm{K}\) and densities exceed \(150\,\mathrm{g\,cm^{-3}}\) \parencite{Vinyoles2017}. The dominant production mechanisms are the proton--proton (\(pp\)) chain and the CNO cycle, generating neutrinos from a few hundred keV up to \(\sim20\,\mathrm{MeV}\) \parencite{Vinyoles2017}.

Unlike reactor or accelerator neutrinos, solar neutrinos are produced deep inside a medium whose density decreases continuously from the core to the photosphere, so that flavour evolution is strongly affected by the MSW effect. For the Large Mixing Angle (LMA) solution favoured by current oscillation data \parencite{NuFit52}, propagation through the Sun is highly adiabatic over most of the relevant energy range: the neutrino remains in the same instantaneous matter eigenstate as the local density changes \parencite{Haxton1986}. Within this approximation, the flavour state produced at radius \(r\) is projected onto the local matter eigenstates, which evolve adiabatically into vacuum mass eigenstates as the neutrino exits the Sun. After propagating over the Earth--Sun distance (1~AU), the mass eigenstates lose mutual coherence due to wave-packet separation and the finite energy resolution of detectors \parencite{Giunti2007}, so the flavour transition probabilities reduce to an incoherent sum,

\begin{equation}
  P_{\nu_e\to\nu_\alpha}^{(f)}(E) =
  \sum_{i=1}^{3}
  |U_{\alpha i}|^2 T_i^{(f)}(E),
\end{equation}

where \(T_i^{(f)}(E)\) is the probability that a neutrino produced through solar reaction \(f\) emerges from the Sun as vacuum mass eigenstate \(\nu_i\). Since production is spatially distributed throughout the solar interior, \(T_i^{(f)}\) is obtained by averaging over the corresponding radial production profile,

\begin{equation}
  T_i^{(f)}(E) = \int_0^{R_\odot} w_f(r)\, P_{\nu_e\to\nu_i}(E,r)\, dr,
\end{equation}

with \(w_f(r)\) the normalised production distribution predicted by the Standard Solar Model (SSM). The solar module in \texttt{tpeanuts} uses radial electron-density profiles together with production distributions for the principal channels (\(pp\), \(pep\), \(hep\), \({}^7\mathrm{Be}\), \({}^8\mathrm{B}\), \({}^{13}\mathrm{N}\), \({}^{15}\mathrm{O}\), \({}^{17}\mathrm{F}\)), using standard solar models such as B16(GS98) and B16(AGSS09met) \parencite{Vinyoles2017}. Optional non-adiabatic corrections based on the Landau--Zener formalism \parencite{Parke1986} are also available for the regions where the adiabatic approximation becomes less accurate. Because solar neutrinos experience both matter-enhanced conversion inside the Sun and, at night, regeneration inside the Earth, they are one of the most powerful probes of neutrino propagation in matter and provide stringent constraints on the solar mixing parameters and on possible non-standard interactions.

\subsubsection{Atmospheric neutrinos}

Atmospheric neutrinos are produced when high-energy cosmic rays interact with nuclei in the atmosphere, generating hadronic cascades dominated by pions and kaons; the decay of these mesons and of the resulting secondary muons produces both muon and electron neutrinos, spanning energies from a few hundred MeV to beyond the TeV scale \parencite{Honda2007}. Unlike solar neutrinos, they are produced over a wide range of altitudes and arrive at terrestrial detectors from all directions, with baselines extending from \(\sim20\,\mathrm{km}\) for downward-going events to nearly the Earth's diameter (\(\sim12\,742\,\mathrm{km}\)) for upward-going neutrinos that cross the planet.

This broad coverage in energy and baseline makes atmospheric neutrinos sensitive to several oscillation regimes at once: downward-going events are dominated by atmospheric or vacuum oscillations, while upward-going neutrinos experience significant matter effects while traversing the Earth's mantle and core. The interplay between the two regimes produces the characteristic oscillation patterns in the \((E,\cos\theta_z)\) plane known as oscillograms. Within \texttt{tpeanuts}, atmospheric-neutrino calculations combine the three-flavour propagation formalism with realistic Earth density profiles, such as the Preliminary Reference Earth Model (PREM) \parencite{Dziewonski1981}, and differential atmospheric flux tables from the Honda parametrisation \parencite{Honda2007} or from cascade solvers such as MCEq \parencite{Fedynitch2015}.

\subsubsection{Reactor antineutrinos}
\label{sec:reactor-source}

Nuclear reactors are intense, well-characterised sources of electron antineutrinos, emitted in the beta decays of the neutron-rich fission fragments of \(^{235}\mathrm{U}\), \(^{238}\mathrm{U}\), \(^{239}\mathrm{Pu}\) and \(^{241}\mathrm{Pu}\), with a spectrum extending up to \(\sim8\,\mathrm{MeV}\) and a rate set by the reactor's thermal power. Because flavour evolution is entirely determined by the vacuum Hamiltonian of \cref{sec:propagation-fundamentals}, and reactor baselines (\(\lesssim2\,\mathrm{km}\) for a near/far experiment such as Daya Bay) traverse negligible Earth matter, reactor antineutrinos follow the exact coherent three-flavour vacuum survival probability, with \(\theta_{13}\) governing \(\bar\nu_e\to\bar\nu_e\) disappearance together with the atmospheric-scale splitting. At this baseline \(\sin^2\Delta_{31}\) and \(\sin^2\Delta_{32}\) are numerically close enough, differing only through the small \(\Delta_{21}\) phase, that what the disappearance channel directly resolves is the effective, ordering-insensitive combination \(\Delta m^2_{ee}=\cos^2\theta_{12}\,\Delta m^2_{31}+\sin^2\theta_{12}\,\Delta m^2_{32}\), rather than \(\Delta m^2_{31}\) itself, which is only recovered after fixing \(\theta_{12}\) and \(\Delta m^2_{21}\) externally (\cref{sec:dayabay-case-study}). \texttt{tpeanuts} implements a full reactor-antineutrino observable chain on top of this probability, from the per-isotope Huber--Mueller emission spectrum \parencite{Huber2011,Mueller2011} through the inverse-beta-decay cross section \parencite{VogelBeacom1999} to a detector-response model folding the true antineutrino spectrum into an observable prompt-energy spectrum, described in \cref{sec:detector-layer} and validated against real Daya Bay data in \cref{sec:dayabay-case-study}.

\subsubsection{Other neutrino sources}

Accelerator-based experiments produce controlled muon-neutrino or antineutrino beams over baselines from a few hundred to more than a thousand kilometres, enabling precision measurements of the atmospheric oscillation parameters, the CP-violating phase and the neutrino mass ordering. As with reactor antineutrinos, incorporating such a source requires only specifying the appropriate initial flavour composition, energy spectrum, baseline and matter-density profile within the same propagation engine, without any change to the oscillation formalism itself.

\subsection{New physics}

Beyond-the-Standard-Model scenarios can be incorporated into the same formalism through modifications of the effective Hamiltonian and, in the case of sterile neutrinos, of the PMNS matrix itself. Two such extensions are considered in \texttt{tpeanuts}: non-standard interactions and sterile neutrinos.

\subsubsection{Non-standard interactions (NSI)}
\label{sec:nsi-formalism}

Although the three-flavour framework introduced above explains the vast majority of neutrino oscillation observations, it is widely regarded as an effective low-energy description rather than a complete theory: the Standard Model does not account for the origin of neutrino masses, nor does it explain the nature of dark matter, the matter--antimatter asymmetry or the hierarchy of fermion masses. Many theories beyond the Standard Model (BSM) therefore predict additional neutrino interactions that would manifest as small deviations from the standard oscillation paradigm.

Among the most studied such scenarios are non-standard interactions (NSI), in which neutrinos experience additional weak-scale interactions with matter mediated by new heavy particles \parencite{Grossman1995}. At energies well below the mediator mass scale, these effects can be described with an effective field theory in which the heavy degrees of freedom are integrated out and encoded in dimensionless effective couplings, providing a model-independent framework for a broad class of BSM scenarios without specifying the underlying microscopic theory \parencite{FarzanTortola2018}.

NSI can in general affect neutrino production, detection and propagation; this work considers only \emph{propagation} NSI, in which new neutral-current-like interactions modify the coherent forward scattering of neutrinos through matter. Since coherent forward scattering only changes the propagation phase, these interactions appear as an additional contribution to the matter Hamiltonian,

\begin{equation}
  H_{\mathrm{mat}} =
  \sqrt{2}G_F n_e
  \left[
  \begin{pmatrix}
    1 & 0 & 0\\
    0 & 0 & 0\\
    0 & 0 & 0
  \end{pmatrix}
  +
  \begin{pmatrix}
    \varepsilon_{ee} & \varepsilon_{e\mu} & \varepsilon_{e\tau}\\
    \varepsilon_{e\mu}^{*} & \varepsilon_{\mu\mu} & \varepsilon_{\mu\tau}\\
    \varepsilon_{e\tau}^{*} & \varepsilon_{\mu\tau}^{*} & \varepsilon_{\tau\tau}
  \end{pmatrix}
  \right],
  \label{eq:nsi-hamiltonian}
\end{equation}

where the dimensionless couplings \(\varepsilon_{\alpha\beta}\), normalised to the electron density \(n_e\) for a given medium composition, quantify the interaction strength relative to the Standard Model weak interaction. Only the traceless part of the diagonal is observable, since a common shift to all three entries is an unobservable global phase; \texttt{tpeanuts} adopts the standard convention \(\varepsilon_{\mu\mu}=0\). The off-diagonal terms introduce flavour-changing coherent interactions with no Standard Model analogue; for antineutrinos, \(\varepsilon_{\alpha\beta}\) is complex-conjugated together with the sign flip of \cref{sec:antinu-convention}. More generally, NSI couple to the individual matter fermions \(e\), \(u\), \(d\) rather than to a single medium, so the coupling above is really \(\varepsilon_{\alpha\beta}(x)=\varepsilon_{\alpha\beta}^{(e+p)}+(n_n(x)/n_e(x))\,\varepsilon_{\alpha\beta}^{n}\), using electrical neutrality (\(n_p=n_e\)) to collapse \(e,u,d\) into two effective pieces that \texttt{tpeanuts} implements separately, \(\varepsilon^{n}\) defaulting to zero to recover \cref{eq:nsi-hamiltonian}; this matters because the Sun and the Earth have different neutron-to-electron ratios, one ingredient of the LMA-Dark degeneracy below. The complete Hamiltonian is \(H=H_{\mathrm{vac}}+H_{\mathrm{mat}}\), with \(H_{\mathrm{vac}}=\frac{1}{2E}U M^2 U^{\dagger}\) unchanged: NSI leave vacuum oscillations unaffected but modify matter propagation, with an impact depending on density, composition and energy, making long-baseline, atmospheric and solar experiments particularly sensitive.

Phenomenologically, NSI modify the effective mixing angles and mass splittings in matter, shift the MSW resonance, alter the day--night regeneration of solar neutrinos, and can introduce degeneracies with the standard oscillation parameters. A well-known example is the LMA-Dark solution \parencite{MirandaTortolaValle2004,Esteban2018}, where a large negative \(\varepsilon_{ee}\) compensates a transformation of the solar mixing angle, yielding oscillation probabilities nearly indistinguishable from the standard Large Mixing Angle (LMA) solution; such degeneracies illustrate the value of combining different classes of neutrino experiments to disentangle standard oscillation physics from possible BSM effects. Current data constrain most individual NSI couplings to be subdominant with respect to the Standard Model weak interaction, but not universally: near degeneracies such as LMA-Dark, large couplings survive precisely because they are compensated by a parameter transformation rather than excluded outright, and couplings involving the \(\tau\) flavour remain only weakly bounded on their own; global analyses combining oscillation, scattering and collider measurements give increasingly stringent, but data-set-dependent, limits \parencite{Biggio2009,Esteban2018}.

From the computational point of view, NSI are a natural extension of the propagation formalism implemented in \texttt{tpeanuts}: since they modify only the matter contribution to the Hamiltonian, the numerical propagation algorithms developed for the standard three-flavour model remain unchanged. The implementation introduces a fully general complex Hermitian NSI matrix, allowing both flavour-conserving and flavour-changing couplings, which enables direct comparison between Standard Model predictions and a broad range of BSM scenarios while preserving the efficiency and flexibility of the underlying PyTorch implementation.

\subsubsection{Sterile neutrinos}
\label{sec:sterile-formalism}

The Standard Model contains three active flavours, \(\nu_e\), \(\nu_\mu\) and \(\nu_\tau\), all of which participate in the weak interaction through charged- and neutral-current processes. Several BSM models, together with a number of experimental anomalies reported over the past decades, have nonetheless motivated the existence of additional neutrino states that are singlets under the Standard Model gauge group. Because they do not couple directly to the weak interaction, these hypothetical states are known as \emph{sterile} neutrinos \parencite{Boser2020}. Although they cannot be produced or detected through Standard Model weak processes, sterile neutrinos may mix with the active flavours through the neutrino mass matrix and thus participate in flavour oscillations as additional propagation eigenstates, modifying the observed oscillation pattern without introducing any new interaction at the detector.

The simplest and most widely studied extension is the \(3+1\) model \parencite{Kopp2013,GiuntiLasserre2019}, in which a single sterile state is added to the three active flavours. The flavour basis becomes \((\nu_e,\nu_\mu,\nu_\tau,\nu_s)\) and the propagation basis four mass eigenstates \((\nu_1,\nu_2,\nu_3,\nu_4)\), enlarging the \(3\times3\) PMNS matrix into a \(4\times4\) unitary mixing matrix, built here as
\begin{equation}
  U_4 = R_{23}\,\Delta\,R_{24}(\theta_{24},\delta_{24})\,R_{34}(\theta_{34},\delta_{34})\,R_{13}\,\Delta^{\dagger}\,R_{12}\,R_{14}(\theta_{14},\delta_{14}),
\end{equation}
extending the three-flavour factorisation of \cref{sec:propagation-fundamentals}, with \(R_{23}\), \(\Delta\), \(R_{13}\), \(R_{12}\) embedded unchanged and three new active-sterile rotations \(R_{14}\), \(R_{24}\), \(R_{34}\), each carrying its own CP phase. This ordering is \texttt{tpeanuts}'s own implementation convention rather than a universal standard, and must be respected exactly: \(R_{13}\) and \(R_{34}\) share the \(\tau\) index, and \(\Delta\) fails to commute with \(R_{34}\) once \(\delta_{\mathrm{CP}}\neq0\), so it is not equivalent to a block-embedded \(U_{\mathrm{PMNS}}\) followed by the sterile rotations. Beyond the three-flavour parameters, the \(3+1\) framework introduces three new mixing angles \(\theta_{14},\theta_{24},\theta_{34}\); \texttt{tpeanuts} parametrises the sterile sector with three raw phases \(\delta_{14},\delta_{24},\delta_{34}\), of which only two are physically independent, consistent with the general \(N\times N\) count of \((N-1)(N-2)/2\) physical phases (three for \(N=4\), including \(\delta_{\mathrm{CP}}\)). The extension also adds a fourth mass eigenstate \(m_4\) and \(\Delta m_{41}^2=m_4^2-m_1^2\), typically \(\gg|\Delta m_{31}^2|\) (of order \(1\,\mathrm{eV}^2\)) in short-baseline-anomaly phenomenology \parencite{Kopp2013}, giving oscillation lengths much shorter than the atmospheric or solar ones.

Unlike NSI, which modify only the matter contribution, sterile neutrinos change the formalism more fundamentally: the vacuum Hamiltonian itself must be enlarged. The formalism of \cref{sec:propagation-fundamentals} is expressed in natural units, \(\hbar=c=1\); \texttt{tpeanuts}'s internal normalisation rescales it by a purely numerical reference length \(L_{\mathrm{scale}}\), together with the explicit \(\hbar c\) conversion needed once \(L_{\mathrm{scale}}\) is expressed in metres rather than natural (inverse-energy) units, without changing the underlying physics. Consistently with the matter potentials below, the enlarged vacuum term is

\begin{equation}
  H_{\mathrm{vac}} = \frac{L_{\mathrm{scale}}}{2E\hbar c}\,U_4\operatorname{diag}\!\left(m_1^2,m_2^2,m_3^2,m_4^2\right)U_4^{\dagger},
\end{equation}

so the new eigenstate contributes its own phase, \(\exp(-im_4^2L/2E)\) in natural units, to the evolution operator. Matter effects are modified too: the sterile state feels neither the charged- nor the neutral-current (NC) potential common to the three active flavours, \(V_{\mathrm{NC}}=\mp\frac{1}{\sqrt2}G_Fn_nL_{\mathrm{scale}}\) (from the neutron density \(n_n\)) — an unobservable global phase in the pure three-flavour case, but a genuine, observable relative potential once this NC-blind sterile state is added: after subtracting the common phase,

\begin{equation}
  H_{\mathrm{mat}} = \operatorname{diag}(V_{\mathrm{CC}},\,0,\,0,\,-V_{\mathrm{NC}}),
\end{equation}

with \(V_{\mathrm{CC}}=\sqrt2 G_F n_e L_{\mathrm{scale}}\). This NC term is an optional addition, active whenever a neutron-density profile is supplied; omitting it recovers the CC-only sterile matter term used as the default here, an approximation once \(\theta_{24}\) or \(\theta_{34}\) are non-zero, and one not to be assumed negligible in general: the NC potential is standard practice for matter-affected active-sterile propagation in the literature, underlying the Earth-matter resonance exploited by atmospheric and long-baseline \(\nu_\mu\)-\(\nu_s\) searches \parencite{Abbasi2020IceCubeSterile}. Either way, the modified matter Hamiltonian alters the effective eigenvalues and can generate additional resonant conversion in dense media.

Phenomenologically, sterile neutrinos may produce anomalous flavour disappearance, unexpected appearance channels, distortions of oscillation spectra and modified matter-induced effects, and have been extensively investigated in connection with the LSND and MiniBooNE anomalies, the reactor antineutrino anomaly and the Gallium anomaly, although no conclusive evidence has yet established their existence \parencite{Boser2020}. Cosmological observations independently constrain light sterile neutrinos through their impact on Big Bang nucleosynthesis, the effective number of relativistic degrees of freedom, and large-scale structure formation.

Within \texttt{tpeanuts}, sterile neutrinos are implemented as a natural extension of the standard three-flavour propagation framework: the oscillation engine has been generalised to support arbitrary Hamiltonian dimensions, so that the \(4\times4\) mixing matrix, the enlarged mass spectrum and the corresponding matter Hamiltonian are treated with the same numerical propagation algorithms developed for the standard model. The current implementation fully supports the \(3+1\) scenario, while its modular design allows straightforward extension to models with multiple sterile states or more general BSM oscillation frameworks.

\subsection{Statistical inference formalism}
\label{sec:inference-formalism}

Comparing an oscillation prediction against real, binned experimental data requires a statistical framework connecting predicted and observed counts, a rule for minimising the resulting statistic with respect to the oscillation and detector parameters, and a way of estimating the uncertainty on the fitted values. \texttt{tpeanuts} implements this framework directly on top of PyTorch's automatic differentiation, so that the same gradient used to build the oscillation Hamiltonian also drives the fit.

For a set of bins with predicted counts \(N_k(\theta)\) and observed counts \(n_k^{\mathrm{obs}}\), the Poisson (Baker--Cousins) likelihood-ratio statistic \parencite{BakerCousins1984}

\begin{equation}
  -2\ln L_{\mathrm{Poisson}} = 2\sum_k\left[N_k(\theta)-n_k^{\mathrm{obs}}+n_k^{\mathrm{obs}}\ln\frac{n_k^{\mathrm{obs}}}{N_k(\theta)}\right]
  \label{eq:poisson-nll-marco}
\end{equation}

is used for absolute event counts, and vanishes exactly when \(N_k=n_k^{\mathrm{obs}}\) bin by bin. For derived observables that are not themselves Poisson-distributed, such as a ratio of counts, a Gaussian statistic with propagated uncertainty is used instead, \(-2\ln L_{\chi^2}=\sum_k[(R_k(\theta)-R_k^{\mathrm{obs}})/\sigma_k]^2\). Externally measured nuisance parameters (background normalisations, energy-scale pulls) are constrained by an additional Gaussian penalty term on \(\theta\) itself, added to either statistic.

The best-fit point \(\hat\theta=\arg\min_\theta[-2\ln L(\theta)]\) is found with the quasi-Newton \texttt{LBFGS} optimiser, using the exact gradient of the total statistic obtained by backpropagating through the full physics chain, oscillation probability, detector response and likelihood, rather than finite differences. The uncertainty on \(\hat\theta\) is estimated through the Laplace approximation, in which the Fisher information at the optimum, \(\mathcal I=\tfrac12\nabla^2_\theta(-2\ln L)\), obtained as the autograd Hessian of the statistic, gives the parameter covariance \(\Sigma=\mathcal I^{-1}\), with \(\sigma_j=\sqrt{\Sigma_{jj}}\). Where the local Laplace approximation is inadequate, in particular along directions the data barely constrain, a direct two-dimensional scan of \(-2\ln L\) over a parameter grid, optionally profiling the remaining free parameters at each grid point, provides confidence contours at Wilks'-theorem asymptotic thresholds, valid under standard regularity conditions but liable to fail near physical boundaries, degeneracies or low statistics; since \texttt{tpeanuts} performs no Monte Carlo (e.g.\ Feldman--Cousins) coverage calibration, these contours should be read as asymptotic, not empirically calibrated.

\subsection{Related computational tools}

\texttt{tpeanuts} interfaces with several public tools, either for cross-validation or as external inputs. \texttt{nuSQuIDS} \parencite{nuSQuIDS}, a general framework for solving neutrino quantum kinetic equations, is used as an independent, actively maintained numerical solver against which vacuum, solar and atmospheric probabilities are cross-validated; PEANUTS remains lighter and more specialised, precisely the property \texttt{tpeanuts} aims to preserve. MCEq \parencite{Fedynitch2015}, which solves the coupled cascade equations of cosmic-ray-induced air showers, and the Honda/HKKM tables \parencite{Honda2007} provide the two atmospheric-neutrino flux inputs; the companion \texttt{CRFlux} library supplies the primary cosmic-ray flux parametrisations (e.g.\ Polygonato) shared by MCEq and the CORSIKA-based sampling pipeline, so that both remain anchored to the same top-of-atmosphere spectrum. \texttt{pymsis} \parencite{Emmert2021} provides empirical atmospheric density/composition profiles as a function of altitude, location and time, replacing constant-density approximations. CORSIKA \parencite{Heck1998} is used in an exploratory capacity to generate primary cosmic-ray samples and shower-level neutrino yields.
